% 解説プリントのコメント（saiten が書く。./tools/kaisetsu.ps1 がこのファイルからプリントを作る）
% 2026-11-22 のスキャン（p4〜p7）の採点から作った
%% matome
\begin{daiji}
  \item 反射・透過の振幅 $r$，$t$，$F$ は複素数である。$\phi^{*}$ を作るときは係数も $r\to r^{*}$ とし，確率流や反射率は $\abs{r}^{2}$，$\abs{t}^{2}$ で書く。$r^{2}$ と $\abs{r}^{2}$ は別物である（問題1 (1)・問題2 (2)・問題3 (2)・問題5 (4)）。
  \item 古典力学との比較は，まず古典の粒子がどうなるかを確定させてから書く。$E>V_{1}$ の段差は古典では必ず越える（$R=0$，$T=1$）。デルタ関数障壁は古典ではどんなエネルギーでも跳ね返し，井戸はどんなエネルギーでも通す（問題1 (4)・問題3 (3)）。
\end{daiji}
\begin{maemo}
  \item 複素数の $\abs{A}^{2}$ を $A^{2}$ と書く：第3章 問題2（2026-10-18・2026-10-21）で同じ書き方をし，第3章 問題5（2026-10-18）では $\psi^{*}$ の中の $R$ を $R^{*}$ にせず実数として扱った。今回は4問で出ている。防ぎ方：$\phi^{*}$ を書く行で係数すべてに ${}^{*}$ を付け，2乗は必ず $\abs{\cdot}^{2}$ で書く。最後に「この係数は実数か」を一度問う。
\end{maemo}
%% mondai 1
\begin{kotaeran}
  \machigai{1}{○}{記述} $k$，$q$ と，干渉の項が打ち消すことは正しい。ただし $\phi^{*}$ を $(\ee^{-\ii kx}+r\ee^{\ii kx})$ と書いて $r^{*}$ にしておらず（$x>0$ も $t\ee^{-\ii qx}$），$j=\hbar k(1-r^{2})/m$，$j=\hbar q t^{2}/m$ となっている。正しくは $\phi^{*}=\ee^{-\ii kx}+r^{*}\ee^{\ii kx}$ で，$j=\hbar k(1-\abs{r}^{2})/m$，$j=\hbar q\abs{t}^{2}/m$。この問題では後で $r$，$t$ が実数になるが，(1) の時点では複素数として扱う。干渉の項は $r^{*}\ee^{2\ii kx}-r\ee^{-2\ii kx}$（$z-z^{*}$ の形で純虚数）で，$\ii k$ を掛けると実数になるので虚部に残らない。
  \machigai{4}{△}{方針} $R=(6-2\sqrt{5})^{2}/16$ と $T=\num{0.854}$ は正しい。$R$ の値は $\num{0.1459}$ なので3桁では $\num{0.146}$（$\num{0.145}$ は切り捨て。$R+T=\num{0.999}$ になっていた）。また「古典粒子では $R=1$，$T=0$」としたが，$E>V_{1}$ の古典の粒子は段差で減速するだけで必ず越えるので $R=0$，$T=1$。量子力学では越えられるエネルギーでも $\num{0.146}$ の割合で反射する，というのが比較の要点である。
\end{kotaeran}
\begin{rikai}
  \item 確率流は $j=(\hbar/m)\Im(\phi^{*}\phi')$。$\phi^{*}$ では係数も複素共役にする。
  \item 古典力学では，$E>V$ の段差・障壁は必ず越え，$E<V$ なら必ず跳ね返る。量子力学ではどちらも $R$，$T$ が0と1の間になる（$E>V$ でも反射し，$E<V$ でもトンネルする）。
  \item 有効数字は次の桁で四捨五入する。$R+T=1$ の検算で丸めの誤りに気づける。
\end{rikai}
%% mondai 2
\begin{kotaeran}
  \machigai{2}{○}{記述} $r=-(\kappa+\ii k)/(\kappa-\ii k)$ と $\abs{r}=1$ は正しい。問われているのは「$x>0$ の確率流が0」だが，問題1 の式 $j=\hbar k(1-r^{2})/m$ に「$r^{2}=\abs{r}^{2}=1$」を入れて $x<0$ の $j=0$ を出している。$r$ は複素数なので $r^{2}\ne\abs{r}^{2}$（$r=\ee^{-2\ii\theta}$ なら $r^{2}=\ee^{-4\ii\theta}$）。$x>0$ では $\phi^{*}\phi'=-\kappa\abs{C}^{2}\ee^{-2\kappa x}$ が実数なので $j=0$，と直接示す（$x<0$ から出すなら，定常状態では $j$ が $x$ によらないことを書く）。
\end{kotaeran}
\begin{rikai}
  \item $x>0$ の $\phi=C\ee^{-\kappa x}$ は「実数の関数×定数」なので $\phi^{*}\phi'$ は実数で，確率流は0。波はしみ込むが，確率は運ばない（しみ込みと透過は別）。
  \item $r^{2}$ と $\abs{r}^{2}$ の区別は問題1 (1) と同じ。
\end{rikai}
%% mondai 3
\begin{kotaeran}
  \machigai{2}{○}{記述} $R=\beta^{2}/(\beta^{2}+k^{2})$，$T=k^{2}/(\beta^{2}+k^{2})$，$T$ が $\beta^{2}$ だけで決まるので $g$ の符号によらないことは正しい。途中で $j_{\mathrm{ref}}=-\hbar k r^{2}/m$，$R=r^{2}=\abs{\beta/(\ii k-\beta)}^{2}$ と書いた。$r=\beta/(\ii k-\beta)$ は複素数なので $r^{2}$ ではなく $\abs{r}^{2}$（最後に絶対値をとったので値は正しい）。問題1 (1) と同じ。
  \machigai{3}{○}{記述} $E\to0$ で $R\to1$，$T\to0$，$E\to\infty$ で $R\to0$，$T\to1$ は正しい。古典との比較が「古典粒子に近づく」「壁が無い状態での古典粒子に近づく」だけで，古典の粒子がどうなるかが書かれていない。デルタ関数障壁（$g>0$）は高さが無限なので古典ではどんなエネルギーでも跳ね返し，井戸（$g<0$）はどんなエネルギーでも通す。量子力学では，高いエネルギーでは障壁でもほとんど透過し，低いエネルギーでは井戸でもほとんど反射する。
\end{kotaeran}
\begin{rikai}
  \item 古典との比較は，(i) 古典ではどうなるか，(ii) 量子ではどうなるか，(iii) どこが違うか，の順に書く。障壁と井戸は，量子では $T$ が同じでも古典では正反対になるので，両方について書く。
\end{rikai}
%% mondai 5
\begin{kotaeran}
  \machigai{4}{○}{記述} $\abs{A}^{2}$ と $\abs{B}^{2}$ を別々に計算して $R$・$T$ を出し，$T$ の式と $R+T=1$ を示したのは正しい。$\abs{A}^{2}=\abs{\cdots}^{2}F^{2}$ と書いたが，$F$ は複素数なので $\abs{F}^{2}$（問題1 (1) と同じ）。また，$R=\abs{B}^{2}/\abs{A}^{2}$，$T=\abs{F}^{2}/\abs{A}^{2}$ とする根拠（入射側と透過側の波数が同じ $k_{1}$ なので速さが等しく，確率流の比が振幅の2乗の比になる）が書かれていない。
\end{kotaeran}
\begin{rikai}
  \item $T=\abs{F}^{2}/\abs{A}^{2}$ が使えるのは両側の波数が同じときだけである。問題1 の段差のように波数が違えば，$T=(q/k)\abs{t}^{2}$ と速さの比が付く。
\end{rikai}
